\section{Introduction}
\label{sec:introduction}

Active Directory Certificate Services (ADCS) issues certificates that can be used for domain authentication through PKINIT~\cite{rfc4556}. Different services record issuance and subsequent use. When a certificate template is changed, used and restored between scans, its current configuration no longer describes the conditions under which an earlier certificate was issued. Investigators must therefore connect directory changes, CA records and KDC authentication outcomes.

SpecterOps documented certificate-service escalation techniques~\cite{schroeder2021certified}. Certificate mapping and the installed Windows updates also affect whether a certificate is accepted for an administrative identity~\cite{microsoft2025strongbinding}. We ask what certificate-specific attribution is justified when identifiers or timing evidence are incomplete, late or contradictory.

\subsection{Certificate Issuance and the Evidence Gap}
In an ADCS deployment, certificate issuance is governed by Active Directory objects known as \textit{Certificate Templates} (stored in the LDAP configuration naming context). When an enterprise template exposes an insecure combination of attributes—such as the \texttt{CT\_FLAG\_ENROLLEE\_SUPPLIES\_SUBJECT} bitflag alongside an Extended Key Usage (EKU) permitting Client Authentication (\texttt{1.3.6.1.5.5.7.3.2}) or Smart Card Logon (\texttt{1.3.6.1.5.2.3.4}), without requiring CA administrator approval—any domain principal granted enrollment permissions can request an X.509 certificate specifying an arbitrary Subject Alternative Name (SAN). The attacker can supply the user principal name (UPN) of a Domain Administrator (e.g., \texttt{Administrator@domain.local}), obtain a cryptographically valid certificate signed by the enterprise root CA, and submit that certificate to the Kerberos Key Distribution Center (KDC) via PKINIT to request a Ticket Granting Ticket (TGT), when the KDC accepts the resulting identity mapping. Certificate-binding policy and patch level can prevent this outcome~\cite{microsoft2025strongbinding}.

Configuration assessment and incident reconstruction answer different questions. Certipy reports template and CA risks~\cite{lyak2021certipy}; a positive finding does not by itself show that a certificate was used with elevated privileges. Whole-network provenance systems connect activity across hosts~\cite{liu2026hades,liu2026cpd}, but a certificate-level explanation still needs an issuer-scoped link between issuance and use.

This study examines three issues. First, a graph classifier can learn template attributes that correlate with generated labels while missing the access-control path. Such shortcut behavior is a known learning problem~\cite{geirhos2020shortcut}; our experiments test it under a stated ADCS permission abstraction. Second, restored configuration attributes cannot identify every intervening issuance. Third, the captured CA and KDC records expose different certificate fields. Existing SAN-based CA/KDC queries provide useful operational correlation~\cite{insecurity2023adcs,splunk2026esc1}. We examine what can still be concluded when the required fields are absent or disagree.

\subsection{Our Contributions}
CertGraph combines graph analysis, certificate-specific evidence reconstruction and budget-constrained remediation. The contributions and their evaluation boundaries are as follows:

\begin{itemize}
\item \textbf{Certificate-specific attribution.} We specify how issuer-scoped issuance records, certificate identifiers and authentication outcomes support possible, definite, conflicting or withdrawn links. Synthetic histories measure the trade-off between retaining candidates and making definite attributions. Controlled Windows matrices and a concurrent transfer workload supply 72 outcome-checked certificate cases. An equal-input SQL comparison tests whether the contract can be expressed relationally.
\item \textbf{Matched analysis of graph learning.} We give a conditional feature-sensitivity result and distinguish the historical feature-channel implementation from a literal bilinear model. The latter reproduces 10,800 held-out predictions on family-disjoint inputs. Its classifications match additive attention, and both models miss the tested long-path positives. The result identifies a verification requirement rather than demonstrating a general advantage for gating.
\item \textbf{Measured remediation trade-offs.} We formulate interdiction with assigned disruption costs and immutable permissions. A metadata-preserving rerun checks 120 graph/budget instances; exhaustive small cases expose gaps from the optimum. A separate path-preserving projection compares the same reachability task. These experiments evaluate the specified instances and costs, not production disruption or a guaranteed optimal policy.
\end{itemize}

\subsection{Paper Organization}
The remainder of this manuscript is organized as follows: Section~\ref{sec:threat_model} formalizes the heterogeneous identity multigraph schema and demonstrates the static snapshot blind spot and ASN.1 SAN concealment on live systems. Section~\ref{sec:certgraph_arch} details the CertGraph GNN architecture and the gating formulation and its implementation boundary. Section~\ref{sec:cipg_provenance} develops the CIPG provenance engine and Invariant~V automaton. Section~\ref{sec:game_theoretic} details the game-theoretic Stackelberg defense and NP-hardness proof. Section~\ref{sec:evaluation} presents empirical results across synthetic and live GOAD cyber-range testbeds. Section~\ref{sec:related_conclusion} reviews related work. Section~\ref{sec:discussion} discusses deployment costs and validity, and Section~\ref{sec:conclusion} concludes the paper.
